\begin{agradecimentos}

A conclusão deste trabalho marca o encerramento de uma etapa construída ao longo de anos de aprendizado, desafios, mudanças de rumo e crescimento pessoal. Embora leve apenas um nome na capa, este documento é resultado da contribuição direta e indireta de muitas pessoas, presentes em conversas, dúvidas tiradas, revisões e incentivos que tornaram possível transformar uma ideia em conhecimento estruturado.

Agradeço à minha orientadora pela disponibilidade, pelos ensinamentos e pela confiança depositada neste trabalho. Sua orientação foi essencial para transformar ideias e questionamentos em uma pesquisa estruturada, conduzida com rigor e responsabilidade.

Agradeço aos professores da graduação em Engenharia Química, que contribuíram para minha formação acadêmica e profissional. Mais do que transmitir conteúdos, eles ajudaram a desenvolver uma forma de pensar baseada na análise crítica, na busca por evidências e na compreensão de sistemas complexos.

À Universidade, agradeço pelas oportunidades de aprendizado, pela infraestrutura oferecida e pelo ambiente que possibilitou meu desenvolvimento ao longo desses anos.

À minha mãe, Dalva, agradeço pelo apoio constante e pelos incentivos recebidos durante toda a graduação, presentes tanto nas decisões mais simples quanto nos momentos de maior incerteza ao longo do curso. A confiança que você depositou em mim, mesmo quando o caminho não estava claro, foi parte do que tornou esta conquista possível.

Ao meu pai, Edson \textit{(in memoriam)}, agradeço por tudo o que me ensinou e por ter sido parte de quem eu me tornei até aqui, mesmo que a pandemia não tenha permitido que ele acompanhasse esta etapa final da minha formação. Esta conquista também é dedicada a ele, com a certeza de que seu lugar nessa trajetória não diminui pela sua ausência física.

À Lavínia, deixo um agradecimento especial por ter compartilhado comigo esta jornada. Obrigada pela paciência nos momentos de maior ansiedade, pelo apoio quando os desafios pareciam maiores do que minha capacidade de enfrentá-los e pela presença constante ao longo desses anos. Seu cuidado esteve presente não apenas nos momentos marcantes, mas também nas pequenas rotinas que sustentam o cotidiano e permitem seguir em frente mesmo diante das dificuldades. Compartilhar a vida com você significou ter ao meu lado alguém que acreditou no meu potencial e me incentivou a continuar. Esta conquista também reflete parte do apoio, da parceria e da construção que realizamos juntas ao longo desse período.

À Stephanie, agradeço pelos anos de amizade construídos durante a graduação. Entre aulas, provas, trabalhos, estágios, preocupações e conquistas, dividimos experiências que tornaram essa caminhada mais leve. Em diversos momentos, quando desistir parecia uma alternativa mais simples do que continuar, seu apoio, suas palavras e sua companhia me ajudaram a seguir adiante, e poucas coisas na graduação foram tão constantes quanto essa amizade.

Ao Raian, agradeço pela amizade construída ao longo dos anos. Sou grata pelo respeito, pelo apoio e pelas conversas compartilhadas em diferentes momentos da minha trajetória, que ficou mais fácil de atravessar com a sua presença e apoio.

Agradeço também aos amigos, colegas e a todas as pessoas que contribuíram direta ou indiretamente para a realização deste trabalho. Cada conversa, cada ajuda oferecida e cada gesto de incentivo tiveram seu valor ao longo dessa caminhada.

Agradeço o apoio financeiro do Programa de Recursos Humanos da Agência Nacional do Petróleo, Gás Natural e Biocombustíveis (PRH/ANP-MCT n\textsuperscript{o}~36/UFBA), suportado com recursos provenientes do investimento de empresas petrolíferas qualificadas na Cláusula de PD\&I da Resolução ANP n\textsuperscript{o}~50/2015, e da Fundação de Amparo à Pesquisa do Estado de São Paulo (FAPESP), processo n\textsuperscript{o}~2024/10318-2.

Os desafios, aprendizados e experiências que fizeram parte desta trajetória reforçaram em mim a convicção de que a engenharia vai além da aplicação de técnicas e modelos: ela é uma ferramenta para compreender sistemas complexos e construir soluções capazes de gerar impacto real.

Que os aprendizados construídos ao longo desta trajetória sirvam de base para continuar transformando problemas complexos em soluções úteis.

\end{agradecimentos}